% ============================================================
%  CHAPTER 3 — Methodology: WBC Framework Design
%  Logical chain: Gap → Proposed Solution → Justification
%  Status: Not started
% ============================================================

\chapter{Methodology}
\label{chap:methodology}

% Opening: restate the identified gap (§2.6) and announce the proposed
% hierarchical framework. Do NOT redefine terms already defined in §2.

\lipsum[14]

% ─── 3.1 FRAMEWORK OVERVIEW ─────────────────────────────────
\section{Framework Overview}
\label{sec:framework-overview}

% Describe the two-layer architecture:
%   Layer 1: Centroidal MPC (planning, ~50–100 Hz)
%   Layer 2: QP-based WBC (execution, ~1 kHz)
% TODO: Insert architecture diagram (figures/framework_overview.pdf)

\begin{figure}[htb]
  \centering
  % \includegraphics[width=0.85\linewidth]{figures/framework_overview.pdf}
  \fbox{\rule{0pt}{4cm}\rule{8cm}{0pt}}  % placeholder — replace with actual figure
  \caption{Overview of the proposed hierarchical WBC framework. The centroidal
           MPC planning layer generates reference trajectories at 1,kHz,
           which are tracked by the QP-based WBC execution layer at 1,kHz.}
  \label{fig:framework-overview}
\end{figure}

\lipsum[15]

\begin{criticalinsight}
  % TODO: Justify the two-layer split.
  % Acknowledge the interface error risk (§6 Controversy 3 / Insight 1).
\end{criticalinsight}

% ─── 3.2 SYSTEM MODEL ───────────────────────────────────────
\section{System Model}
\label{sec:system-model}

\subsection{Floating-Base Dynamics}
\label{sec:floating-base}

% CANONICAL NOTATION — defined here for the whole document.
% All subsequent chapters MUST reuse these symbols without redefinition.

The humanoid robot is modeled as a floating-base rigid multibody system with
$n$ actuated joints. The generalized configuration is $\q = [\q_b\transpose,
\q_j\transpose]\transpose \in \mathbb{R}^{6+n}$, where $\q_b \in SE(3)$ represents
the unactuated base pose and $\q_j \in \mathbb{R}^{n}$ the joint coordinates.
The equations of motion are:

\begin{equation}
  \M \ddq + \Cg = \bm{S}\transpose \torque + \Jc\transpose \fc
  \label{eq:eom}
\end{equation}

\noindent where $\M \in \mathbb{R}^{(6+n)\times(6+n)}$ is the joint-space inertia matrix,
$\Cg \in \mathbb{R}^{6+n}$ collects Coriolis, centrifugal and gravitational terms,
$\bm{S} = [\bm{0}_{n\times 6}, \bm{I}_n]$ is the actuation selection matrix,
$\torque \in \mathbb{R}^n$ is the joint torque vector, $\Jc$ is the contact
Jacobian, and $\fc \in \mathbb{R}^{3k}$ is the contact wrench ($k$ contact points).

\lipsum[16]

\begin{criticalinsight}
  % TODO: Discuss floating-base unactuability and its consequences for WBC.
\end{criticalinsight}

\subsection{Centroidal Dynamics}
\label{sec:centroidal}

% Reference: Orin & Goswami (2013) — centroidal is sufficient for planning.
% Notation: \lmom (linear momentum), \amom (angular momentum), \com (CoM position).

\begin{equation}
  \begin{bmatrix} \dot{\lmom} \\ \dot{\amom} \end{bmatrix}
  = \sum_{i} \begin{bmatrix} \bm{f}_{c,i} \\ \bm{r}_i \times \bm{f}_{c,i} \end{bmatrix}
  + \begin{bmatrix} m\bm{g} \\ \bm{0} \end{bmatrix}
  \label{eq:centroidal}
\end{equation}

\lipsum[17]

% ─── 3.3 QP-BASED WBC FORMULATION ───────────────────────────
\section{QP-Based WBC Formulation}
\label{sec:qp-wbc}

% Formalize the QP: decision variables, cost, constraints.
% Decision vector: \zqp = [\ddq\transpose, \torque\transpose, \fc\transpose]\transpose

\begin{equation}
  \min_{\ddq,\, \torque,\, \fc} \;
    \sum_{i} w_i \norm{\J_i \ddq + \dot{\J}_i \dq - \ddxtask_i^{\mathrm{ref}}}_{\bm{W}_i}^2
  \label{eq:qp-cost}
\end{equation}

\noindent subject to:
\begin{align}
  \M\ddq + \Cg &= \bm{S}\transpose\torque + \Jc\transpose\fc
    & &\text{(floating-base dynamics)} \label{eq:qp-dyn} \\
  \Jc\ddq + \dot{\Jc}\dq &= \bm{0}
    & &\text{(rigid contact)} \label{eq:qp-contact} \\
  \torque_{\min} &\leq \torque \leq \torque_{\max}
    & &\text{(torque limits)} \label{eq:qp-torque} \\
  \fc &\in \mathcal{FC}
    & &\text{(friction cone)} \label{eq:qp-friction}
\end{align}

% TODO: Discuss friction cone linearization (inner polyhedral approximation).
% TODO: Discuss task hierarchy and weight tuning.

\lipsum[18]

\begin{criticalinsight}
  % TODO: Address the QP infeasibility risk at contact transitions (§8.4 Risk).
  % Soft constraint relaxation strategy must be described here.
\end{criticalinsight}

% ─── 3.4 CENTROIDAL MPC PLANNING LAYER ──────────────────────
\section{Centroidal MPC Planning Layer}
\label{sec:mpc-planning}

% Reference: Caron (2020), Winkler/TOWR (2018)
% Describe receding-horizon formulation, footstep planning, interface to WBC.

\lipsum[19]

\begin{criticalinsight}
  % TODO: Quantify the interface error between MPC output and WBC input.
\end{criticalinsight}

% ─── 3.5 CONTACT MODELING ───────────────────────────────────
\section{Contact Modeling and Mode Transitions}
\label{sec:contact-modeling}

% Key controversy (§6 Problem 1): rigid LCP vs. soft contact.
% MUST justify choice here — dedicate explicit sub-section.
% Tie to simulator choice (PQ2 — to be resolved).

\subsection{Rigid Contact Model}
\lipsum[20]

\subsection{Soft Contact Model}
\lipsum[21]

\subsection{Contact Mode Transition Strategy}
% Contact mode switching introduces hybrid dynamics discontinuities.
% Describe how the WBC handles stance $\leftrightarrow$ swing transitions.
\lipsum[22]

\begin{criticalinsight}
  % TODO: Explicitly justify the contact model choice.
  % Recommend sensitivity analysis if time permits (§18 reviewer anticipation).
\end{criticalinsight}

% ─── 3.6 STATE ESTIMATION ───────────────────────────────────
\section{State Estimation}
\label{sec:state-estimation}

% EKF on IMU + joint encoders for base pose/velocity (§3.3 of memory).
% Key trade-off: speed vs. accuracy of base velocity — critical for WBC stability.

\lipsum[23]

\begin{criticalinsight}
  % TODO: Discuss propagation of base velocity estimation error into the WBC.
\end{criticalinsight}

% ─── 3.7 TASK HIERARCHY AND PRIORITY ASSIGNMENT ─────────────
\section{Task Hierarchy and Priority Assignment}
\label{sec:task-hierarchy}

% Define task stack: e.g., (1) contact stability, (2) CoM tracking,
% (3) swing foot tracking, (4) end-effector / manipulation task,
% (5) posture regularization.

\lipsum[24]

\begin{criticalinsight}
  % TODO: Explain how manipulation tasks are integrated without triggering
  % task-priority conflicts with locomotion tasks (§7.2 Methodological Gap).
\end{criticalinsight}